% Generated draft: requires mathematical and source review.
\documentclass[11pt]{article}
\usepackage[margin=25mm]{geometry}
\usepackage{fontspec}
\setmainfont{Latin Modern Roman}
\usepackage{amsmath,amssymb,mathtools,unicode-math}
\setmathfont{Latin Modern Math}
\usepackage{enumitem,needspace,xurl,hyperref}
\setlength{\parindent}{0pt}
\setlength{\parskip}{5pt}
\setlength{\emergencystretch}{3em}
\newcommand{\N}{\mathbb{N}}
\newcommand{\Z}{\mathbb{Z}}
\newcommand{\Q}{\mathbb{Q}}
\newcommand{\R}{\mathbb{R}}
\newcommand{\C}{\mathbb{C}}
\newcommand{\F}{\mathbb{F}}
\newcommand{\A}{\mathbb{A}}
\newcommand{\PP}{\mathbb{P}}
\newcommand{\E}{\mathbb{E}}
\newcommand{\eps}{\varepsilon}
\newcommand{\dd}{\,\mathrm{d}}
\begin{document}
\section*{UnsolvedMath ID: 800783}
\textbf{Source problem code:} catalogue entry
\subsection{Definitions and mathematical statement}
Let \(\operatorname{End}(\C,0)\) denote the holomorphic germs \(f:(\C,0)\to(\C,0)\) fixing the origin. The multiplier of such a germ is \(\lambda=f'(0)\). A germ is holomorphically linearizable at the origin if there exists a locally invertible holomorphic germ \(h\), with \(h(0)=0\) and \(h'(0)\ne0\), such that
\[
h\circ f\circ h^{-1}(z)=\lambda z.
\]
For the irrationally indifferent interpretation of the source's term ``elliptic,'' write \(\lambda=e^{2\pi i\theta}\), where \(\theta\in\R\setminus\Q\). Let \(q_n\) be the denominators of the continued-fraction convergents of the representative of \(\theta\) modulo \(\Z\) in \((0,1)\), with \(q_0=1\). The angle \(\theta\) satisfies the Brjuno condition if
\[
\sum_{n=0}^{\infty}\frac{\log q_{n+1}}{q_n}<\infty.
\]
The source's phrase ``\(\lambda\) is a Brjuno number'' is interpreted as this condition on its rotation angle.

The problem asks for families
\[
f_{\lambda,a}(z)=\lambda z+z^2g_a(z),\qquad a\in M,
\]
where each \(g_a\) is holomorphic near the origin, having the property that, for every admitted parameter pair \((\lambda,a)\) with \(\lambda=e^{2\pi i\theta}\) irrationally indifferent,
\[
f_{\lambda,a}\text{ is holomorphically linearizable at }0
\quad\Longleftrightarrow\quad
\sum_{n=0}^{\infty}\frac{\log q_{n+1}}{q_n}<\infty.
\]
The source does not specify the parameter space \(M\), any regularity of dependence on \(a\), or whether all irrationally indifferent multipliers must occur in each sought family. These choices remain unresolved.

The source asks, in particular, whether every elliptic polynomial with a fixed point of multiplier \(\lambda\) is holomorphically linearizable at that fixed point if and only if \(\lambda\) satisfies the Brjuno condition. After translating the fixed point to the origin, this concerns the local germ of a polynomial \(P\) with \(P(0)=0\) and \(P'(0)=\lambda\). It separately asks the same universal question for polynomials of degree exactly three. The general polynomial question does not state a degree restriction; the intended inclusion or exclusion of degree-one polynomials requires clarification.

As background, the supplied record attributes a dichotomy to Ilyashenko (1979, reference label [I1]) and Perez-Marco (2001, reference label [P9]). In its supplied wording, for every \(g\in\operatorname{End}(\C,0)\) and every \(\lambda\in\C^*\) that is not a root of unity, set
\[
F_{\lambda,a}(z)=\lambda z+azg(z),\qquad a\in\C.
\]
The stated alternatives are that the origin is a Siegel point for every \(a\in\C\), or that it is a Cremer point for every \(a\in\C\setminus K\), where \(K\subset\subset\C\) is a bounded exceptional set of logarithmic capacity zero, and hence of planar Lebesgue measure zero. Here a Siegel point means an irrationally indifferent fixed point whose germ is holomorphically linearizable, and a Cremer point means an irrationally indifferent fixed point whose germ is not holomorphically linearizable. With these meanings, the supplied hypothesis \(\lambda\in\C^*\) not a root of unity is insufficient to place the multiplier in the irrationally indifferent domain: the additional condition \(|\lambda|=1\) is absent. The dichotomy is therefore retained as attributed source background with an unresolved hypothesis, rather than asserted here as a verified theorem in its supplied generality.
\subsection{Short English statement}
Find families of holomorphic germs fixing the origin for which local holomorphic linearizability is equivalent to the Brjuno condition on the irrational rotation multiplier. The source also asks whether this equivalence holds for every elliptic polynomial, and specifically for every elliptic cubic polynomial. Family parameters, the general polynomial degree range, and a hypothesis in the attributed background dichotomy require clarification.
\subsection{Sources}
\begin{itemize}
\item Local UnsolvedMath registry, permanent ID \(800783\), source problem code \texttt{catalogue entry}, title ``Local Discrete Holomorphic Dynamics --- Open Problem OP1.'' The supplied record attributes the problem to Abate, with proposed year 2011, and labels its status ``open.'' This status is supplied registry metadata, not an independently verified assessment.
\item Abate, \emph{Open Problems in Local Discrete Holomorphic Dynamics} (2011), Problem OP1. Supplied source URL: \url{https://arxiv.org/abs/1107.4876}. The supplied provenance records access on 2026-10-16 and extraction by \texttt{arxiv\_latex}; these are source-provided provenance statements, not a claim of independent access or verification here.
\item The supplied background attributes its dichotomy to Ilyashenko, 1979, reference label [I1], and Perez-Marco, 2001, reference label [P9]. Full bibliographic entries and separate links for these two references are not supplied, and their statements and hypotheses have not been independently verified here.
\end{itemize}
\end{document}
